\documentclass[10pt]{article}
\usepackage[margin=20mm]{geometry}
\usepackage{amsmath,amssymb,mathtools,fontspec,unicode-math,xurl,hyperref}
\IfFontExistsTF{DejaVu Serif}{\setmainfont{DejaVu Serif}}{\setmainfont{Latin Modern Roman}}
\setmathfont{latinmodern-math.otf}
\setlength{\parindent}{0pt}\setlength{\parskip}{6pt}\sloppy
\title{Technical review of Diagonal cycles and anticyclotomic Iwasawa theory of modular forms}
\author{AI-assisted technical review}
\date{}
\begin{document}
\maketitle
Abstract: This review examines the hypotheses, local conditions, norm-relation dependencies, and selected arithmetic applications of the bound accepted-author version. The weight boundaries and the distinction between nonzero Selmer rank and nonvanishing of a distinguished class are consistent in the inspected arguments. No substantive defect is established. The geometric construction of the cycles, the regulator identities, and the general Euler-system machinery have not been independently reconstructed in full.\par
\section*{Identity and scope}
The reviewed work is by Francesc Castella and Kim Tuan Do, Diagonal cycles and anticyclotomic Iwasawa theory of modular forms, arXiv:2303.06751v2 [paper]. The arXiv version date is 30 September 2025; the PDF itself is dated 2 October 2025. Both dates are retained. The source is identified as the accepted-author version, not as a byte-verified copy of the publisher version for DOI 10.4171/jems/1735 (https://doi.org/10.4171/jems/1735). The frozen 2025 identifier is unchanged.\par
All references to pages below use one-based PDF pagination, agreeing with printed pages. Actual text reading covered pages 1–7, 10, 21–22, 24–35, 45–48, and 54–57. The full 60-page PDF was mechanically parsed, which does not constitute scientific reading. The review concentrates on Theorems A–C, 3.4.2, 4.2.1, 4.3.1, 4.3.3, and the deductions in Sections 6.5–6.7 and 7.5–7.7. Earlier cycle calculations and Sections 5, 6.1–6.4, and 7.1–7.4 remain substantial dependencies.\par
The basic representation is \[V_{f,\chi}=V_f^\vee(1-r)\otimes\chi^{-1},\] with \(f\) of weight \(2r\geq2\), \(p\nmid6N_f\), \(p\) split in \(K\), and \(f\) ordinary at the selected prime. The assumptions on the conductor of \(\chi\), coprimality, and \(p\nmid h_K\) are retained. The last condition identifies the bottom \(p\)-primary ring-class field \(K[1]\) with \(K\). The definite applications additionally require \(N^-\) squarefree with an odd number of prime factors, residual absolute irreducibility and \(p\)-distinguishedness, \(p>2r-2\), and non-CM \(f\). The indefinite applications in Section 7 impose the classical Heegner hypothesis; the manuscript explicitly does not establish them for every indefinite factorization of the level.\par
\section*{Local conditions and quantifiers}
\subsection*{Weight boundaries}
Lemma 4.1.2, page 31, uses the convention that the cyclotomic character has Hodge–Tate weight \(-1\). The weights of the ordinary subrepresentation and quotient are \[\begin{array}{c|cc}
 & V^+ & V^-\\ \hline
\mathfrak p & -j-r & -j-1+r\\
\overline{\mathfrak p} & j-r & j-1+r .
\end{array}\] For \(0\leq j<r\), the two \(V^+\) weights are negative and the two quotient weights are nonnegative. At the endpoint \(j=r-1\), the quotient weight at \(\mathfrak p\) is zero, so it remains on the nonnegative side. For \(j\geq r\), both weights at \(\mathfrak p\) are negative, and both at \(\overline{\mathfrak p}\) are nonnegative; at \(j=r\) the latter subrepresentation has weight zero. These checks agree with \[\operatorname{Sel}_{\rm BK}
=\begin{cases}
\operatorname{Sel}_{\rm ord,ord},&0\leq j<r,\\
\operatorname{Sel}_{\rm rel,str},&j\geq r.
\end{cases}\] The equality uses the cited local comparison results, not the weight table alone. Their de Rham and exceptional-invariant hypotheses have not been independently re-proved here. A weight-zero endpoint by itself is not a demonstrated counterexample.\par
Theorem 4.2.1 separates the two projected diagonal classes. The first lies in the relaxed–strict condition; the conjugate projection lies in the ordinary–ordinary condition. Equations (4.3)–(4.5) identify the corresponding summands of the balanced local filtration. The argument at primes away from \(p\) uses purity and conjugate self-duality to obtain vanishing of the rational local cohomology; the lattice local conditions are the propagated ones. That is a different argument from imposing integral unramifiedness without propagation.\par
\subsection*{Rational and integral image hypotheses}
Theorem 4.3.1 bounds a finite-level Selmer group by one dimension conditional on a non-torsion bottom class. Its proof lists absolute irreducibility and two required Galois elements, then invokes non-CM image results. Theorem 4.3.3 separately assumes big image and lists the residual and integral analogues. Consequently, a conclusion over \(\Lambda^-_K\) is not obtained merely by silently replacing the integral hypotheses with the rational ones.\par
The stated characteristic-ideal conclusions are containments in the displayed direction, with squared Euler-system indices. The later analytic main-conjecture divisibilities invert \(p\), and in the indefinite case extend coefficients to the specified unramified completion. They are not asserted as integral equalities of characteristic ideals. The source cites the general Jetchev–Nekovář–Skinner machinery and explicitly notes on page 33 that it was not publicly available at the time of writing, pointing also to an exposition. That dependency is recorded rather than represented as an independently verified theorem in this review.\par
\section*{Norm relations and arithmetic deductions}
\subsection*{Construction boundary}
Pages 21–22 explain how the quotient construction removes an extra \(\ell-1\) factor before Rubin’s norm-relation adjustment. The compactly supported cohomology argument uses that every Künneth contribution in total degree at most two for a product of three affine curves has a degree-zero factor. This supports the displayed comparison at the inspected level; the full preceding correspondence calculations have not been redone.\par
The Euler factor is attached to the dual Tate-twisted representation and geometric Frobenius. Pages 24–29 pass to CM Hida families, impose both the central-character condition (sd) and the residual distinction condition (dist), and identify the two character products \(\xi_1\xi_2\) and \(\xi_1\xi_2^c\). The general character is obtained by twisting a constructed special case, rather than by assuming the special case has arbitrary infinity type. The original modular-curve correspondences, their integral compatibilities, and the complete regulator construction remain outside the reconstructed portion. Bars and duals in extracted formulas were not promoted to typographical findings without image confirmation.\par
\subsection*{Rank zero}
For Theorem 6.5.1, the relevant range is \(0\leq j<r\): in the opposite range the sign already forces the central value to vanish. The chosen auxiliary character gives a unit factor in the triple-product factorization. The explicit reciprocity law then implies that the bottom Euler-system class has nonzero image in \[H^1(K_{\mathfrak p},V_{f,\chi})/
H^1_{\rm ord}(K_{\mathfrak p},V_{f,\chi}).\] The Euler-system theorem first makes \(\operatorname{Sel}_{\rm rel,str}\) a one-dimensional space generated by that class.\par
The two global duality sequences on pages 46–47 then have the required linear-algebra consequence. The map from this one-dimensional Selmer group to the one-dimensional local quotient is nonzero and therefore an isomorphism. Its kernel is \(\operatorname{Sel}_{\rm ord,str}=0\). The first sequence forces \(\operatorname{Sel}_{\rm rel,ord}\) to have dimension one; its inclusion of \(\operatorname{Sel}_{\rm rel,str}\) is then an equality. The nonzero local image in the second sequence forces \(\operatorname{Sel}_{\rm ord,ord}=0\). Lemma 4.1.2 identifies this group with the Bloch–Kato group. The argument is thus not an inference from a nonzero class directly to a zero Selmer group with identical local conditions.\par
The indefinite rank-zero argument uses the other projected class and the opposite weight range \(j\geq r\). The reciprocity law in (7.14)–(7.15) gives a nonzero localization at \(\overline{\mathfrak p}\). The manuscript applies the corresponding duality argument and local comparison. The auxiliary-character nonvanishing theorems, local dimensions, and reciprocity identities are identified inputs; this inspection does not supply new proofs of them.\par
\subsection*{Rank one and specialization}
Theorem 6.7.1 first obtains a non-torsion Iwasawa class from a nonzero \(p\)-adic \(L\)-function. It then uses the injection (6.19) to prove positive rank of the finite-level Selmer group. These are distinct from claiming that this particular class has nonzero specialization.\par
For clarity, the algebra underlying the rank inference is sound for a finitely generated module \(M\) of positive rank over \(\mathcal O[[T]]\). Localizing at the height-one prime \((T)\) gives a module of positive free rank over a discrete valuation ring. Hence \(M/TM\), after inverting the uniformizer of \(\mathcal O\), has positive dimension. It does not follow that a selected element of \(M\) survives modulo \(T\). The manuscript preserves exactly this distinction: equality of the finite-level dimension to one is conditional on nonvanishing of the distinguished specialized class. The same distinction is present in Theorem 7.7.1 and (7.19). The control injections themselves, including the propagated local conditions, were checked as named dependencies, not independently established in full.\par
\section*{Assessment and remaining limits}
No substantive error is established in the inspected statements and reductions. The assessment is No Issue Found within the inspected scope. In particular, this does not certify the geometric cycle construction, all tame correspondence identities, the analytic interpolation formulas, or the external Euler-system machinery.\par
The acquired metadata, exact-version arXiv record, publisher record, author publication list, and bounded title/identifier searches did not identify a separate supplement, correction to the bound version, or computational repository. The author publication list warns that its files may differ from arXiv; no unverified replacement was substituted. Absence of auxiliary material is not proved, and no author code was run.\par
Item | State | Scope\par
Issue judgment | No Issue Found | Declared hypotheses and selected proof reductions\par
Repair/adoption | Not applicable | No established defect\par
Scientific decision | Pending | No human or independent acceptance\par
Remaining work | Visual inspection | Source evidence pages and every report page\par
\section*{Supporting checks}
> 00000TOKEN.\par
Paper ID: `M25-0356`.\par
Reasoning and writing are by this assistant session. Native tools perform file checks, extraction, actual compilation, parsing, and rendering; they are not independent reviewers. Controller UI observations are attributed to the controller, and backend equivalence is unverified. Image reads have failed with `EMFILE`; no rendered image is credited as viewed. Delivery completion remains false until actual hash-bound inspection is available.\par
9 F. Castella and K. T. Do, Diagonal cycles and anticyclotomic Iwasawa theory of modular forms, arXiv:2303.06751v2 (30 September 2025; internal PDF date 2 October 2025), https://arxiv.org/abs/2303.06751v2. EMS Press, DOI 10.4171/jems/1735, https://ems.press/journals/jems/articles/14299326, consulted 27 October 2026. F. Castella, publication list, https://web.math.ucsb.edu/\textasciitilde{}castella/, consulted 27 October 2026.\par
\end{document}
